% !TeX root = ../papers/nc-p3.tex
\section{二十次二重曲線での二次曲面の射影による反対根の正規形}
\label{sec:quadric-kernel-projection}
\begin{theorem}\label{thm:quadric-kernel-projection}
定理\ref{thm:cofactor-sectional-kernel}で $d=20$ かつ核像が曲面である場合、
八次曲面の正規化はその核像の既約二次曲面 $Q$ である。入力への有限双有理射
$\nu:Q\to\mathbb P^3$ は基点のない四空間 $V\subset H^0(Q,\mathcal O_Q(2))$ による。
\[
 M=\operatorname{coker}N=\nu_*\mathcal O_Q(1),\qquad
 V\otimes H^0(Q,\mathcal O_Q(1))\xrightarrow{\ \sim\ } H^0(Q,\mathcal O_Q(3)).
\]
滑らかな二次曲面でも二次錐でも成立する。各印は入力の長さ三繊維であり、反対の印はその繊維の張る平面を表す。
二根行列の入力基底変更・出力合同による同値の必要十分条件は、偏極二次曲面の同型が一方の $V$ を他方に運ぶことである。
\end{theorem}
\begin{proof}
核グラフ解消での $H,L$ は $H^2=8$, $HL=4$, $L^2=2$ である。Hodge指数の等号により
$H\equiv2L$ なので両射が潰す曲線は一致する。入力の正規化と出力二次曲面はともに正規で、
両方へ固有双有理な解消の繊維は連結である。各射が他方の全繊維を一点に写すので両正規曲面は同型である
（両射の像グラフが各正規曲面へ有限双有理となる）。正規化上の入力束は $\mathcal O_Q(2)$ に等しい。
滑らかなら二母線の次数を比べ、錐なら最小解消 $\mathbb F_2$ 上への引戻しを比べればよい。
両解消のPicard群では数値同値と線形同値は一致する。入力座標は従って所要の四空間を与える。

核点の四座標の評価により $M\to\nu_*\mathcal O_Q(1)$ がある。
$M$ は整超曲面上の階数一極大Cohen--Macaulay層で一般同型なので単射。
表示から $\chi(M(t))=4(t+1)^2$、右辺でも二次曲面のHilbert多項式で
$\chi(\mathcal O_Q(2t+1))=(2t+2)^2$。よって同型である。次数一の表示に関係がないので積写像は同型。
各閉点で生成元数は有限繊維の長さとなり、四つの核座標の繊維への評価が全射する。
階数一点なら長さ三の繊維を消す一次式は一つで、その係数が $N(p)$ の像直線である。
\end{proof}

\begin{proof}[四空間からの行列の再構成と同値]
四空間とその基底を固定し、$U=H^0(Q,\mathcal O_Q(1))$ とする。積写像
\[
  \operatorname{Sym}^2 V\otimes U\longrightarrow H^0(Q,\mathcal O_Q(5))
\]
の核の基底を四本の二次式列に並べる。実際の窓ではこれが $N$ の表示の関係四空間そのものである。
列の定数変換で対称となる可逆変換を線形方程式で求める。
より一般に基点なし有限双有理なデータで上の一次積が同型なら、この手順は二次式四列の極小表示を与える。
実際 $M'=\nu_*\mathcal O_Q(1)$ は $H^1(M'(t))=0$ であり、さらに $H^2(M'(-1))=0$ なので
$1$-正則である。次数一の積の全射により次数零の四元だけで切断加群を生成する。
負次数切断は零、中間コホモロジー消滅により加群はCohen--Macaulayで、射影次元一。
Hilbert分子 $4-4z^2$ が関係の次数と本数を与える。
得た行列について対称化、十成分の独立性、階数一および二の軌跡を直接検査できる。

同じ基底の生成元に対する二つの対称表示が $N,N'=NP$ なら $P^tN=NP$ であり、
$P^t$ は $M'$ の大域自己準同型を与える。その環は定数だけである。
有限双有理射上での正規化の可逆層の自己準同型は分数体の乗法で、
正規化上の各局所環を保つものに限られるからである。
$H^0(M')=U$ 上で見ると $P$ は定数スカラーである。従って
入力と $U$ の基底を除き根は一意に復元され、逆向きには核像と入力への射が根から復元された。
ここで得た同値条件は標識した根行列についてであり、四つの矢空間の順序を含む代数の分類には
中間座標と全次数延長がさらに必要である。
\end{proof}
\path{computations/verify_quadric_kernel_projection.py} は滑らかな二次曲面の四つの $(2,2)$ 式について
一次積、四関係、対称化の一意性、十成分の独立性を有理数上で検算する。
一次積が全射なら基点はなく、正次数の豊富な束を用いるため射は有限である。
表示の行列が像の一点で階数三となる検算は一般繊維長一も保証する。
